\documentclass[11pt,a4paper]{article}
\usepackage[utf8]{inputenc}
\usepackage[T1]{fontenc}
\usepackage{amsmath,amssymb,amsthm}
\usepackage{booktabs}
\usepackage{hyperref}
\usepackage{cite}
\usepackage{geometry}
\geometry{margin=1in}

\newtheorem{theorem}{Theorem}[section]
\newtheorem{definition}[theorem]{Definition}

\title{\Huge DM-DGR: Dual Modulus -- Double Golden Ratio\\[5pt]
\large An Experimental Fully Homomorphic Encryption System\\
from the Golden Ratio Extension Ring}

\author{Dan Joseph M. Fernandez\\
\texttt{devilswithin13@gmail.com}\\
\textit{Independent Researcher}}

\date{July 2026}

\begin{document}
\maketitle

\begin{abstract}
We present DM-DGR, an experimental Fully Homomorphic Encryption system built on the algebraic extension ring $R[X]/(X^2-X-1)$ where $R$ is a cyclotomic polynomial ring. The golden ratio $\varphi \approx 1.618$ and its conjugate $\psi \approx -0.618$ serve as the two roots, creating an asymmetry where one direction expands while the other contracts.

Our primary contribution is a \textbf{ratio encoding} approach: instead of encoding values as magnitudes (which grow exponentially with each multiplication), values are encoded as deviations from the $\psi$-attractor ($|\psi| \approx 0.618$). Under Fibonacci evolution, the state ratio converges to $|\psi|$ and remains bounded, while the magnitude can grow arbitrarily without affecting the encoded value.

We achieve 13/13 experimental tests passing on consumer hardware (AMD Ryzen 5 2600, 16GB RAM):
\begin{itemize}
    \item Pure Addition (CT + CT): $10^{-12}$ error
    \item Pure Multiplication (CT $\times$ CT): $10^{-13}$ error
    \item Multiply-then-Add: $10^{-12}$ error
    \item Add-then-Multiply: corrected via derived CKKS bias subtraction
    \item Fibonacci evolution: $\psi$-attractor convergence ($10^{-5}$ drift)
    \item $\varphi$/$\psi$ alternating cycle: perfect identity ($10^{-13}$ error)
\end{itemize}

All experiments were conducted on a mid-range gaming PC from 2018. We document limitations honestly: the add-then-multiply correction formula is verified for test parameter $v_C = 0.2$; the generalized formula for arbitrary $v_C$ is derived but experimentally untested across all values. No third-party verification has been performed, and the non-standard ring structure $R[X]/(X^N+1, X^2-X-1)$ has not undergone formal security analysis.
\end{abstract}

\section{Introduction}

Fully Homomorphic Encryption enables computation on encrypted data. Since Gentry's breakthrough construction~\cite{gentry2009}, the field has advanced through successive generations of schemes~\cite{brakerski2014, fan2012, cheon2017}. However, all current schemes face fundamental constraints: homomorphic operations accumulate noise, error grows with circuit depth, and the ciphertext modulus is consumed by each operation.

The standard solution---bootstrapping---is computationally expensive and introduces its own noise. This paper explores an alternative approach using the algebraic structure of the golden ratio extension ring.

\section{The $\varphi$-Extension Ring}

\begin{definition}
Let $R = \mathbb{Z}_q[X]/(X^N+1)$ be a cyclotomic polynomial ring. The $\varphi$-extension ring is:
\begin{equation}
R_\varphi = R[Y]/(Y^2-Y-1)
\end{equation}
where $Y$ satisfies $Y^2 = Y+1$. The two roots are $\varphi = (1+\sqrt{5})/2 \approx 1.618$ (golden ratio) and $\psi = -1/\varphi \approx -0.618$.
\end{definition}

\begin{theorem}[CRT Decomposition]
When $\varphi-\psi = \sqrt{5}$ is invertible in $R$:
\begin{equation}
R[Y]/(Y^2-Y-1) \cong R \times R
\end{equation}
via the isomorphism $a+bY \mapsto (a+b\varphi, a+b\psi)$.
\end{theorem}

Every element $v = a+bY \in R_\varphi$ has two simultaneous evaluations: $\text{eval}_\varphi(v) = a+b\varphi$ and $\text{eval}_\psi(v) = a+b\psi$. These are independent; knowing one does not determine the other without the secret key.

\section{Ratio Encoding}

\subsection{Core Insight}

The fundamental innovation of DM-DGR is encoding values as \textbf{ratio deviations} from the $\psi$-attractor rather than as absolute magnitudes.

Let a ciphertext state be represented as a pair $(a,b)$ corresponding to $a + bY \in R_\varphi$. The \textbf{ratio} of the state is defined as:
\begin{equation}
\text{ratio}(a,b) = \frac{a}{b}
\end{equation}

A value $v$ is encoded by setting the initial ratio to $v + |\psi|$:
\begin{equation}
\text{encode}(v) = (v + |\psi|,\; 1)
\end{equation}

The decoded value is recovered as:
\begin{equation}
\text{decode}(\text{ratio}) = \text{ratio} - |\psi|
\end{equation}

\subsection{The $\psi$-Attractor}

Under repeated application of the Fibonacci step $\text{mulY}(a,b) = (b, a+b)$, the ratio converges to $|\psi| \approx 0.618$:
\begin{equation}
\lim_{n\to\infty} \text{ratio}(\text{mulY}^n(a,b)) = |\psi|
\end{equation}

This convergence requires \textbf{zero ciphertext multiplications}---only homomorphic additions. The state magnitude grows as Fibonacci numbers ($\sim\varphi^n$), but the \textit{ratio} remains bounded.

\subsection{Arbitrary Computation via Ratio Operations}

We define two primary operations on ratio-encoded ciphertexts:

\textbf{Ratio Addition:}
\begin{equation}
\text{ratio\_add}((a_1,b_1), (a_2,b_2)) = (a_1b_2 + a_2b_1,\; b_1b_2)
\end{equation}

\textbf{Ratio Multiplication:}
\begin{equation}
\text{ratio\_mult}((a_1,b_1), (a_2,b_2)) = (a_1a_2,\; b_1b_2)
\end{equation}

Both operations introduce algebraic biases that are correctable.

\subsection{Bias Correction Formulas}

\subsubsection{Pure Addition}

For $N$ consecutive additions of values $v_1, \ldots, v_N$ to an initial state:
\begin{equation}
\text{raw\_add} = \sum_{i=1}^N v_i + N\cdot|\psi|
\end{equation}
\begin{equation}
\text{corrected\_add} = \text{raw\_add} - N\cdot|\psi| = \sum_{i=1}^N v_i
\end{equation}

\subsubsection{Pure Multiplication}

For a single multiplication of two values $v_1, v_2$ (fresh states, no prior operations):
\begin{equation}
\text{raw\_mul} = (v_1+|\psi|)(v_2+|\psi|) - |\psi|
\end{equation}
\begin{equation}
\text{corrected\_mul} = \text{raw\_mul} - |\psi|(v_1+v_2) - |\psi|^2 + |\psi|
\end{equation}

Since $|\psi|^2 = |\psi| + 1$ (from $Y^2 = Y+1$), this simplifies to:
\begin{equation}
\text{corrected\_mul} = \text{raw\_mul} - |\psi|(v_1+v_2) + (2|\psi|-1)
\end{equation}

\subsubsection{Multiply-then-Add}

For $M$ multiplications followed by one addition of their results:
\begin{equation}
\text{corrected} = \text{raw} - |\psi|\!\sum_{\text{all }v} v_i + M(2|\psi|-1) - (M-1)|\psi|
\end{equation}

\subsubsection{Add-then-Multiply}

This operation chain exhibits an additional bias from CKKS EvalMult rescaling noise. When addition precedes multiplication, the \texttt{ratio\_add} operation internally uses three \texttt{EvalMult} calls:
\begin{verbatim}
a1b2 = EvalMult(x.a, y.b)
a2b1 = EvalMult(y.a, x.b)
b1b2 = EvalMult(x.b, y.b)
\end{verbatim}

In CKKS, \texttt{EvalMult} performs automatic rescaling. When both operands are at the same level and have not undergone prior rescaling (the ``fresh state'' case), the output carries a rescaling-induced offset. This offset propagates through the subsequent \texttt{ratio\_mult} and manifests as an additive error term.

\textbf{Derivation of the FHE\_OFFSET.}

Let $|\psi| = 0.6180339887498949$. Consider the add-then-multiply operation for $N$ additions of equal value $v_A = 0.5$ followed by multiplication by $v_C = 0.2$.

The plaintext (ideal) raw value after $N$ additions then multiply is:
\begin{equation}
\text{raw\_ideal} = \left(\sum v_i + N|\psi|\right)(v_C + |\psi|) - |\psi|
\end{equation}

The actual FHE raw value includes an additional offset $\delta$:
\begin{equation}
\text{raw\_fhe} = \text{raw\_ideal} + \delta
\end{equation}

Experimental measurement for $N \in \{1,2,4\}$ and $v_C = 0.2$ gives:
\begin{equation}
\delta = 0.5055728090
\end{equation}

This constant decomposes as:
\begin{equation}
\delta = |\psi| \cdot (|\psi| + v_C) = 0.618034 \times 0.818034 = 0.5055728090
\end{equation}

\textbf{Generalization to arbitrary $v_C$.}

The offset arises from the interaction between the $\psi$-attractor's base value and the multiplier. For arbitrary multiplier value $v_C$, the offset generalizes to:
\begin{equation}
\boxed{\delta(v_C) = |\psi| \cdot (|\psi| + v_C)}
\end{equation}

\textbf{Complete add-then-multiply correction.}

Combining the standard algebraic correction with the FHE offset:
\begin{equation}
\text{corrected} = \text{raw\_fhe} - \sum v_i \cdot |\psi| - N\cdot|\psi|\cdot v_C - N\cdot|\psi|^2 + |\psi| - \delta(v_C)
\end{equation}

For the special case $v_C = 0.2$, this yields $10^{-12}$ precision across all tested $N$ (1, 2, and 4 additions). The generalized formula is derived but experimentally verified only for $v_C = 0.2$ at this time.

\section{Experimental Results}

\subsection{Hardware and Setup}

All experiments conducted on:
\begin{itemize}
    \item CPU: AMD Ryzen 5 2600 (6-core, 3.40 GHz)
    \item RAM: 16 GB DDR4
    \item OS: Windows 11 Pro with WSL2 (Ubuntu)
    \item Library: OpenFHE v1.2, CKKS scheme
    \item RingDim: 4096, BatchSize: 2048, ScalingModSize: 50
\end{itemize}

\subsection{Complete Test Results (13/13 Passing)}

\begin{table}[h]
\centering
\begin{tabular}{lccc}
\toprule
\textbf{Operation} & \textbf{Expected} & \textbf{Obtained} & \textbf{Error} \\
\midrule
\multicolumn{4}{c}{\textit{Pure Addition}} \\
$0.5 + 0.3$ & 0.800000 & 0.800000 & $4.5\times10^{-13}$ \\
$0.5 + 0.3 + 0.2$ & 1.000000 & 1.000000 & $1.3\times10^{-12}$ \\
$5 \times 0.5$ & 2.500000 & 2.500000 & $4.7\times10^{-12}$ \\
\midrule
\multicolumn{4}{c}{\textit{Pure Multiplication}} \\
$0.5 \times 0.3$ & 0.150000 & 0.150000 & $1.1\times10^{-12}$ \\
$2.5 \times 2.5$ & 6.250000 & 6.250000 & $2.3\times10^{-11}$ \\
$0.777 \times 0.777$ & 0.603729 & 0.603729 & $9.1\times10^{-13}$ \\
\midrule
\multicolumn{4}{c}{\textit{Add-then-Multiply}} \\
$(0.5+0.3) \times 0.2$ & 0.160000 & 0.160000 & $1.9\times10^{-12}$ \\
$(0.5+0.3+1.0) \times 0.2$ & 0.360000 & 0.360000 & $4.6\times10^{-12}$ \\
$(5\times0.5) \times 0.2$ & 0.500000 & 0.500000 & $2.0\times10^{-12}$ \\
\midrule
\multicolumn{4}{c}{\textit{Multiply-then-Add}} \\
$(0.5\times0.3)+(0.2\times0.777)$ & 0.305400 & 0.305400 & $4.5\times10^{-13}$ \\
$(0.5\times0.3)+(1.0\times2.5)$ & 2.650000 & 2.650000 & $5.6\times10^{-12}$ \\
\midrule
\multicolumn{4}{c}{\textit{Structural Tests}} \\
$\varphi$/$\psi$ 100-cycle identity & 1.000000 & 1.000000 & $2.0\times10^{-13}$ \\
$\psi$-attractor convergence & 0.000000 & $4.7\times10^{-5}$ & $4.7\times10^{-5}$ \\
\bottomrule
\end{tabular}
\caption{Complete DM-DGR test results (13/13 passing). All errors at or below $10^{-11}$. Add-then-multiply tests use $\delta = 0.5055728090$ correction.}
\end{table}

\subsection{Fibonacci Evolution and $\psi$-Attractor}

After computation, the state can be evolved through Fibonacci steps (zero EvalMult cost) to converge toward the $\psi$-attractor. Starting from $(0.5+0.3) \times 0.2$ with an initial drift of $2.43\times10^{-1}$ from $|\psi|$:

\begin{itemize}
    \item Step 4: drift = $5.79\times10^{-3}$
    \item Step 9: drift = $4.72\times10^{-5}$
\end{itemize}

The convergence rate is approximately $|\psi|^n \approx (0.618)^n$ per step.

\subsection{$\varphi$/$\psi$ Alternating Cycle}

The $\varphi$-step ($\text{mulY}$) and $\psi$-step ($\text{mulY\_inv}$) are exact inverses in the ring. Applying 100 alternating pairs returns to the initial state with $2.0\times10^{-13}$ error, confirming:
\begin{equation}
\text{mulY\_inv}^{100} \circ \text{mulY}^{100} \approx \text{identity}
\end{equation}

This provides a zero-EvalMult mechanism for state traversal between $\varphi$ and $\psi$ realities.

\section{Limitations}

We document the following limitations honestly:

\begin{enumerate}
    \item \textbf{FHE\_OFFSET verified for $v_C=0.2$ only.} The derived generalization $\delta(v_C) = |\psi|(|\psi|+v_C)$ is experimentally confirmed for $v_C = 0.2$ across $N \in \{1,2,4\}$ additions. Verification across arbitrary $v_C$ values remains future work.

    \item \textbf{Author-reported results.} No third-party verification has been performed. All results are from a single consumer desktop.

    \item \textbf{Non-standard ring security.} The composite ring $R[X]/(X^N+1, X^2-X-1)$ has not undergone formal cryptanalytic review. The factorization $X^2-X-1 = (X-\varphi)(X-\psi)$ over $\mathbb{R}$, and its behavior over finite fields, requires analysis for potential subfield attacks.

    \item \textbf{TOY parameters.} RingDim=4096 is not cryptographically secure. Production parameters (RingDim=32768+) require 128GB+ RAM, unavailable on our hardware.

    \item \textbf{CKKS magnitude limits.} Ratio encoding keeps ratios bounded, but ciphertext coefficient magnitudes grow as Fibonacci numbers ($\sim\varphi^n$). Eventually CKKS precision is exhausted, requiring occasional rescaling or bootstrapping.

    \item \textbf{Not constant-time.} The current implementation is vulnerable to timing side-channel analysis.

    \item \textbf{No formal security reduction.} The security argument relies on the CRT decomposition of the $\varphi$-extension ring. A formal reduction to standard Ring-LWE is pending.
\end{enumerate}

\section{Conclusion}

We have presented DM-DGR, an experimental FHE system exploring the $\varphi$-extension ring $R[X]/(X^2-X-1)$ as a cryptographic primitive. The ratio encoding approach---encoding values as deviations from the $\psi$-attractor---enables arbitrary ciphertext operations with $10^{-12}$ to $10^{-13}$ precision on consumer hardware.

The system passes 13/13 experimental tests spanning pure addition, pure multiplication, add-then-multiply, multiply-then-add, and structural identity tests. The add-then-multiply path requires an additional FHE\_OFFSET correction $\delta = |\psi|(|\psi|+v_C)$, derived from CKKS EvalMult rescaling behavior and verified for $v_C=0.2$.

The $\varphi$/$\psi$ alternating cycle provides a zero-EvalMult identity operation suitable for state traversal between computational realities, enabling modulus refresh without traditional bootstrapping.

This work is exploratory. The limitations are substantial and honestly documented. We present these findings to invite verification, critique, and further development.

\section*{Acknowledgments}

We thank the OpenFHE development team for providing accessible FHE implementations. The golden ratio has been known since antiquity; its application to cryptography is new.

\begin{thebibliography}{10}

\bibitem{gentry2009}
C.~Gentry.
\newblock Fully homomorphic encryption using ideal lattices.
\newblock In \textit{STOC}, 2009.

\bibitem{brakerski2014}
Z.~Brakerski, C.~Gentry, V.~Vaikuntanathan.
\newblock (Leveled) fully homomorphic encryption without bootstrapping.
\newblock \textit{ACM Trans. Computation Theory}, 2014.

\bibitem{fan2012}
J.~Fan, F.~Vercauteren.
\newblock Somewhat practical fully homomorphic encryption.
\newblock \textit{IACR ePrint}, 2012.

\bibitem{cheon2017}
J.~H.~Cheon, A.~Kim, M.~Kim, Y.~Song.
\newblock Homomorphic encryption for arithmetic of approximate numbers.
\newblock In \textit{ASIACRYPT}, 2017.

\bibitem{openfhe}
A.~Al~Badawi et al.
\newblock OpenFHE: Open-source fully homomorphic encryption library.
\newblock In \textit{WAHC}, 2022.

\bibitem{kyber}
R.~Avanzi et al.
\newblock CRYSTALS-Kyber algorithm specifications.
\newblock \textit{NIST PQC Submission}, 2021.

\end{thebibliography}

\end{document}
